\chapter{War Communism}

The hostility of the outside world was only one of the hazards confronting the Bolsheviks after their assumption of power. The revolution in Petrograd had been bloodless; but sharp fighting occurred in Moscow between Bolshevik units and military cadets loyal to the Provisional Government. Dispossessed political parties began to organize against the authority of the Soviets. Communications were halted by a strike of railway workers, whose trade union was controlled by Mensheviks. Administrative services were disrupted; and hooligans took advantage of the anarchic conditions to riot and plunder. Six weeks after the revolution a governmental decree brought into being the All-Russian Extraordinary Commission (Cheka) for ``combating counter-revolution and sabotage''; and local Soviets were invited to set up similar commissions. A few days later a revolutionary tribunal was set up to try ``those who organize risings against the Workers' and Peasants' Government, who actively oppose or do not obey it, or who incite others to oppose or disobey it''. It was not till June 1918 that the revolutionary tribunal pronounced its first death sentence. But indiscriminate killings both of Bolsheviks and of their adversaries occurred in many parts of the country; and the Cheka became increasingly busy rounding up active opponents of the regime. 

In April 1918 several hundred anarchists were arrested in Moscow; in July the Cheka was called on to suppress an attempted coup by SRs, who assassinated the German Ambassador---apparently in protest against the Brest-Litovsk treaty. During the summer of 1918 two prominent Bolshevik leaders were assassinated in Petrograd, and shots were fired at Lenin in Moscow. The ferocity with which the civil war was fought heightened the tension. Atrocities on one side were matched by reprisals on the other. ``Red terror'' and ``White terror'' both entered the political vocabulary. 
% 注：修复了多栏/跨行乱序 ("...called (Brest-Litovsk treaty. During the summer of 1918 two prominent on to suppress an attempted coup... Bolshevik leaders were assassinated...") 以及清除了 "j.on", "IIterror", "//ary.^" 等乱码。

These desperate conditions were reflected in the total disarray of the economy. During the war, production had been crippled and distorted by military needs, and by the absence of agricultural and industrial workers at the front. The revolution itself, and the ravages of the civil war, completed the picture of economic, social and financial disintegration; hunger and cold overtook large sectors of the population. Initial Bolshevik remedies for economic ills did not go beyond the proclamation of such general principles as equal distribution, nationalization of industry and of the land, and workers' control. In the first months of the revolution many industrial enterprises were taken over, sometimes by state organs responsible to the Supreme Council of National Economy (Vesenkha), sometimes by the workers themselves. For agriculture the Bolsheviks, who still had little power in the countryside, had adopted the programme of the SRs, and proclaimed the ``socialization'' of the land and its equal distribution among those who tilled it. What happened in fact was that the peasants seized, and distributed among themselves, the estates, large and small, of the land-owning nobility, and the holdings of well-to-do peasants, commonly dubbed kulaks, who had been enabled to accumulate land by the Stolypin reforms. None of these measures arrested the decline of production. In finance, the banks were nationalized and foreign debts repudiated. But it was impossible to collect regular taxes or to frame a state budget; current needs were met by resort to the printing-press. For six months the regime lived from hand to mouth. 
% 注：去除了 "h These", "l|of", "bf economic", "(who" 等边缘识别乱码。

Then the gathering storms of the civil war and the economic collapse drove the government in the summer of 1918 to the more drastic policies later known by the ambiguous name of ``war communism''. Food was the first priority. The workers in towns and factories were hungry. In May the word went out to organize ``food detachments'' to visit the countryside and collect grain from kulaks and speculators---the ``rural bourgeoisie''---who were believed to be hoarding it. A decree of June 11, 1918, provided for the creation in the villages of ``committees of poor peasants'', which ``under the general direction of the People's Commissariat of Supply (Narkomprod)'' were to supervise the collection, distribution and despatch to the towns of grain and other agricultural produce. Lenin hailed the constitution of these committees as ``the October, i.e. proletarian, revolution'' in the countryside, and thought that it signalized the transition from the bourgeois to the socialist revolution. But the experiment was short-lived. The decree, like others of the period, was easier to write than to enforce. The spontaneous action of the peasants in the first year of the revolution resulted in the division of the land among a multiplicity of small cultivators living at subsistence level---an increase in the number, and reduction in the size, of units of cultivation which contributed nothing either to the efficiency of agriculture or to the supply of food to the towns, since the small producer was more likely to consume what he produced for his own needs. Poor peasants were not easily organized; and rivalry sprung up between the committees and the village Soviets. Class stratification in the villages was real enough. But the criteria of classification of the peasantry as kulaks, middle and poor peasants, were uncertain and fluctuating, and were partly dictated by the political requirements of the moment. Kulak, in particular, became a term of abuse directed by party propaganda against peasants who incurred the wrath of the authorities through failure to comply with demands for the delivery of grain. Nor could the poor peasants be counted on, as party leaders in Moscow expected, to act as allies of the government against the kulaks. The poor peasant was conscious of the oppression which he suffered at the hands of the kulak. But his dread of the state and its minions was often greater; and he was apt to prefer the evil that he knew to the menace of a remote authority. 
% 注：合并了 "Narkom- prod" 和 "des- jpatch"。修复了 "June II" 为 "June 11"。去除了 "I were", "between 1 the", "ixTld." 等乱码。

In December 1918 the committees of poor peasants were abolished, and the authorities switched their appeal to the so-called ``middle peasants'', who rose above the indigent level of the ``poor peasants'', but did not qualify for the label of ``rich peasants'' or ``kulaks''. But in the chaos of the civil war no expedient could stimulate agricultural production. The authorities from time to time invoked the cherished socialist goal of large-scale collective cultivation. A number of agricultural communes or ``collective farms'' (Kolkhozy) were founded by communist idealists, some of them foreigners, on the basis of working and living in common. But they scarcely contributed to the problem of feeding the cities. ``Soviet farms'' (Sovkhozy) were set up by the Soviet Government, by provincial or local Soviets, or sometimes by industrial enterprises under the control of Vesenkha, for the specific purpose of providing food for famished urban and factory workers; they employed wage-labour, and were sometimes spoken of as ``socialist grain factories''. But they made little headway against the resistance of the peasants, who saw in the Sovkhozy a return to the large landed estates broken up by the revolution, especially when, as often happened, they were established on confiscated estates, and employed the managers taken over from the old regime. Lenin on one occasion repeated a saying alleged to have been current among the peasants: ``We are Bolsheviks, but not communists; we are for the Bolsheviks because they drove out the landlords, but we are not for the communists because they are against individual holdings.'' 
% 注：去除了 "I label", "ci\dl war" (修复为 civil war), "produc- tion", "' were founded", "Govern- ment", "T \ T T workers" 等识别错误。

In industry, war communism may be said to have begun with a decree of June 28, 1918, nationalizing every important category of industry. This seems to have been inspired, partly by the growing menace of the civil war, partly by the desire to forestall spontaneous seizures of factories by workers without the knowledge or authority of Vesenkha---what a writer of the period called ``elemental-chaotic proletarian nationalization from below''. But formal nationalization was of little account. What mattered was to organize and administer what had been taken over---a function which workers' control had proved unable to exercise. This was the task of Vesenkha, which set up a number of ``centres'' or ``chief committees (glavki)'' to manage whole industries; some industrial undertakings were administered by local authorities. Chaotic conditions called urgently for centralized control, which may, however, sometimes have aggravated the chaos. Few of the qualifications and skills required for industrial production were available to the new regime. Industry at all levels continued to be run in practice by those who had worked in it before the revolution, and who now manned the ``centres'' and ``glavki''. Party members were sometimes assigned to top positions, but lacked the experience to make themselves effective. Senior directors, managers and engineers, whose services were quickly recognized as indispensable, were known as ``specialists'', and were rewarded with higher salaries and privileges. Industrial production was, however, increasingly dominated by the emergencies of the civil war. The demands of the Red Army were paramount. Effort had to be concentrated on a few essential industries at the expense of the rest. Small-scale enterprises employing only a handful of workers, and artisan industry both in the towns and in the countryside, were mainly immune from controls, but were frequently hampered by lack of materials. Manpower was mobilized for the front. Transport broke down. Supplies of raw materials were exhausted, and could not be replenished. Of the many statistics illustrating the catastrophic decline of industry perhaps the most revealing were those which recorded the depopulation of the big cities. In the three years after the revolution Moscow lost 44.5 per cent of its population, Petrograd, where the industrial concentration was heaviest, 57.5 per cent. The Red Army took its toll of the able-bodied; and masses of people drifted away to the countryside where, if anywhere, food might still be found. 
% 注：去除了 "Ifrom below", "{glavki)", "[urgently", "Wacked" (修正为 lacked), "' directors" 等错误，并将数字 "44-5"、"57 5" 恢复为正常的 "44.5"、"57.5"。

The problems of distribution were no less recalcitrant. The aim announced in the party programme of replacing private trade by ``a planned system of distribution of commodities on an all-state scale'' was a remote ideal. A decree of April 1918 authorizing Narkomprod to acquire stocks of consumer goods for exchange against peasant stocks of grain remained a dead letter. Plans to enforce rationing and fixed prices in the towns broke down in face of the shortage of supplies and the absence of any efficient administration. Trade flowed, where it flowed at all, in illicit channels. Traders, sufficiently numerous to acquire the familiar nickname of ``bagmen'', travelled round the country with supplies of simple consumer goods which they exchanged with peasants for foodstuffs to be sold at exorbitant prices in the towns. ``Bagmen'' were frequently denounced by the authorities, and threatened with arrest or shooting, but continued to prosper. Some attempt was made to use the existing machinery of the cooperatives, and control was established, not without friction, over the central cooperative organs. Since money was rapidly losing its value, schemes were hatched for the barter of commodities between town and country; but the goods wanted by the peasant were also in short supply. In the critical year of the civil war, when the survival of the regime seemed to hang by a thread, and the territory even nominally controlled by it was being constantly contracted by inroads of the White armies, the method by which the essential needs of the Red Army, of the factories engaged in war production, and of the urban population were met was the crude method of requisitioning, dictated and justified by military necessity. To keep the Red Army supplied was the over-riding task of economic policy, and little attention could be spared for civilian needs or civilian susceptibilities. It was above all the widespread requisitioning of surpluses of grain which led the peasants, once the danger from the Whites was over, to rebel against the harshnesses of war communism. 
% 注：修正了 "famiUar", "friction,,", 并去除了 "|||peasant" 和 "fTTTTf ■>'>1'"。

War communism had important consequences for the organization of labour. The initial hope that, while compulsion would have to be applied to landlords and members of the bourgeoisie, the labour of the workers would be regulated by voluntary self-discipline was soon frustrated. ``Workers' control'' over production, exercised in every factory by an elected factory committee, which had been encouraged in the first flush of revolution, and had played a role in the take-over of power, soon became a recipe for anarchy. In the rapidly thickening crisis atmosphere of January 1918 Lenin significantly quoted the familiar, ``He who does not work, neither shall he eat'', as ``the practical creed of socialism''; and the People's Commissar for Labour spoke of ``sabotage'' and necessary measures of compulsion. Lenin had a good word to say for piece-rates and for ``Taylorism''---a fashionable American system for improving the efficiency of labour, which he himself had once denounced as ``the enslavement of the man to the machine''. Later he supported a campaign for the introduction of what was called ``one-man management'' in industry---the direct antithesis of ``workers' control''. The party congress of March 1918 which voted to ratify the Brest-Litovsk treaty also demanded ``draconian measures to raise the self-discipline and discipline of workers and peasants''. These proposals, like the Brest-Litovsk treaty itself, excited the indignation of the then Left opposition, in which Bukharin and Radek played leading parts. 
% 注：清除了段尾的 "^e'^" 和 "*L"。

The revolution had spot-lighted the ambiguous role of the trade union in a workers' state. Relations between Soviets of Workers' Deputies and trade unions, both purporting to represent the interests of the workers, had been a crux since the earliest days of the revolution, when the strongest unions were dominated by the Mensheviks. When the first All-Russian Congress of Trade Unions met in January 1918, the Bolsheviks had secured a majority, though the Mensheviks and other parties were also well represented. The congress had no difficulty in calling the factory committees to order on the ground that the particular interest of a small group of workers must yield to the general interest of the proletariat as a whole. Only a few anarchist delegates opposed the decision to convert the committees into organs of the unions. Here too the principle of the centralization of the authority dispersed by the revolution was already at work. 

The issue of the relation of the trade unions to the state was far more stubbornly contested. Were the unions to be an integral part of the apparatus of the workers' state like other Soviet institutions? Or would they retain the function of defending specific interests of the workers independently of other elements of the workers' state? The Mensheviks, and some Bolsheviks, arguing that, since the revolution had not yet outlived its bourgeois-democratic stage, the unions still had their traditional role to play, stood out for complete independence of the unions from the state. But Zinoviev, who presided, had no difficulty in securing a comfortable majority for the official Bolshevik view that, in the process of the revolution, the trade unions must ``inevitably be transformed into organs of the socialist state'', and in that capacity must ``undertake the chief burden of organizing production''. Declining production, and the needs of a desperate situation, made this mandate vital. To raise labour productivity, to improve labour discipline, to regulate wages and to prevent strikes were responsibilities which the trade unions, in partnership with Vesenkha and other state organs, were now required to assume. The distinction between the functions of the trade unions and those of the People's Commissariat of Labour (Narkomtrud) became mainly formal; most of the principal officials of Narkomtrud were henceforth trade union nominees. 
% 注：修复了极其严重的多栏错序（"arguing that, since the revolution had i|Of the revolution, the trade unions must “inevitably be trans- not yet outlived its bourgeois-democratic stage..."）并去除了大量乱码如 "ljunions", "UNarkomtrud)", "'Ipfficials", "^"。

The emergency of the civil war revived and kept alive the mood of enthusiasm generated by the revolution itself, and made strict measures of discipline acceptable. In April 1919, with the civil war now at its height, general conscription for military service was ordered; and this soon came in practice to include the drafting of labour for essential work. About the same time labour camps were instituted for offenders sentenced to this form of punishment by the Cheka or by the ordinary courts, who were to be employed on work at the direction of Soviet institutions. The severest category of these camps, known as ``concentration camps'', was reserved for those who had engaged in counter-revolutionary activities in the civil war, and who were to be detailed for particularly arduous work. But appeals could also be made to voluntary self-discipline. In May 1919 Lenin called on the workers for what were named ``Communist Saturdays'', when some thousands of workers in Moscow and Petrograd volunteered to work over-time without pay to speed up the despatch of troops and supplies to the front; and this precedent was followed a year later. The institution of udarniki, or shock workers, to carry out specially important work at high speed dated from this time. Without this combination of harsh compulsion and spontaneous enthusiasm the civil war could not have been won. 
% 注：清除了 "-^^Ithe", "■^Ifand", "7^" 等乱码。

Early in 1920, with the defeat of Denikin and Kolchak, the military emergency had been overcome. But it made way for the equally grave problems of almost total economic collapse; and it seemed logical that these problems should be met by the same forms of discipline which had brought victory in the field. Trotsky as People's Commissar for War made himself the champion of the conscription and ``militarization'' of labour to pave the way for economic revival. During the period of war communism the trade unions had been brushed aside. Labour had been conscripted for work behind the military fronts; and, when fighting ceased, military units were converted into ``labour battalions'' for necessary work of reconstruction. The first ``revolutionary army of labour'' was formed in the Urals in January 1920. Now, however, that the civil war was over, the mood changed. Those who from the first had looked askance at measures of compulsion applied to the workers, those who stood for the independence of the trade unions, and those who for other reasons resented Trotsky's pre-eminence in the party, joined in attacking his masterful procedures. He defended his policies in face of mounting opposition at the party congress in March 1920, and secured Lenin's support. The outbreak of the Polish war stilled the voices of dissent. But when the war ended in the autumn of 1920, and the last embers of the civil war had been stamped out in the south, fierce opposition arose in the party to the continued conscription of labour and the virtual by-passing of the trade unions. Trotsky, impressed by the vast and urgent problems of economic reconstruction, and irritated by trade union resistance to his plans, added fuel to the flames by demanding a ``shake-up'' of the unions. Lenin parted company with Trotsky on the issue; and a bitter debate of unprecedented dimensions raged throughout the winter---only to be resolved when the policies of war communism were finally abandoned at the party congress of March 1921. 
% 注：去除了 "(askance"，修复了 "pre- eminence" 和 "proce- dures" 断字。

Party attitudes to war communism were divided and ambivalent. The conglomeration of practical policies collectively known by that name was approved as necessary and proper by all but a small minority of dissidents. But interpretations of its character diverged widely---more widely, perhaps, in retrospect than at the time. The first eight months of Soviet rule had broken the power of the landlords and the bourgeoisie, but had not yet brought into being a socialist economic order. In May 1918 Lenin still spoke of an ``intention . . . to realise the transition to socialism''. The sudden introduction in the summer, under the name of war communism, of measures which seemed to many Bolsheviks a foretaste of the future socialist economy, was treated by more prudent party members simply as a forced response to an emergency, an abandonment of the cautious advance hitherto pursued, a plunge---necessary, no doubt, but rash and full of hazards---into uncharted waters. This view gained in popularity when the civil war ended, and the burdens of war communism seemed no longer tolerable; and it became the accepted line when peasant revolt finally forced a decision to abandon war communism in favour of NEP. 
% 注：清除了 "'lAUi", 修复了 "toler- able" 和 "ambi- valent"。

Other communists, on the other hand, hailed the achievements of war communism as an economic triumph, an advance into socialism and communism more rapid than had hitherto been deemed possible, but none the less impressive on that account. Industry was comprehensively nationalized; and, if industrial production still declined, Bukharin could write complacently of ``the revolutionary disintegration of industry'' as ``a historically necessary stage''. The progressive devaluation of the ruble could be described as a blow struck against bourgeois capitalists, and a prelude to the moneyless communist society of the future, when everything would be shared according to needs. Already, it was claimed, the market had been largely eliminated as the agency of distribution. Grain surpluses were requisitioned from the peasants; and the main foodstuffs were in principle rationed to the urban population. Industry worked chiefly on government orders. Labour was organized and allocated in response, not to the dictates of the market, but to social and military needs. After the civil war the realities of a desperate economic situation clashed too obviously with this Utopian picture to make it seriously tenable. But many party consciences were troubled by its abandonment; and the divergences of opinion about the character of war communism repeated themselves as divergences about the character and permanence of NEP.
% 注：清除了 "(of", "^ society", "(||largely", "[Iwere", "I were", "11 and", "(but", "^ ^'/.realities", "l^pwith", "i_ . land" 等大量结尾段落的杂乱标记。